\scanpage{25}\mspage{24}
Finalement, on travaillera avec la top. définie d'abord (base formée avec les ``intervalles''), et on appellera néanmoins par abus de langage, les $[a,b]$, $\mathcal L_{\leq a}$, $\mathcal L_{\geq b}$ (ainsi que $\mathcal L$) \uline{sous-tronçons} \unsure{ordonnés|ordinaires} (ou \unsure{biordonnés|bien donnés|bien disons|binominés}) ``fermés'', [NB on les suppose non réduits à un pt, i.e.
\[
\deleted{\unsure{[|i}} a<b\quad \text{resp.}\quad a\ne\mathrm{or}(\mathcal L)\quad \text{resp.}\quad \deleted{\unsure{a}}b\ne\mathrm{ex}(\mathcal L)\;]
\]

Un \uline{sous-tronçon \inserted{fermé} large} est une \deleted{sous-tronçon} partie qui est soit un ss-tronçon, soit \unsure{réduit|réduite} à un pt.

Pour une telle partie, $T$, on
% Widely spaced strokes cross the whole lower block, from pose to the last example.
% The close zigzags and local horizontal cancellations show an earlier level of correction.
\deleted{pose}
\[
\begin{aligned}
\deleted{\partial T=}&\deleted{\{t\in T\mid t\in \overline{\mathcal L\setminus T}\}}\quad\text{\deleted{bord}}\\
 &\text{\deleted{\unsure{frontière dans la top.|frontière dans la top}}}\\
 &\text{\deleted{absolu de $T$ $=$ ensemble}}\\
 &\text{\deleted{canonique}}\\
 \deleted{={}\unsure{T\cap\mathcal L|T\cap\overline{\mathcal L}}}&\quad\text{\deleted{fermé de ses extrémités}}\\
 &\deleted{(\mathrm{NB}\ \mathrm{card}\,\partial T\in\{0,1,2\})}
\end{aligned}
\]

\deleted{On voit que pour un tronçon fermé (\uline{NB} non réduit à un pt) on a}

\[
\begin{aligned}
 \deleted{\partial[a,b]=\{a,b\}}&\qquad \deleted{\partial \mathcal L_{\leq a}=\{a\},\quad \partial \mathcal L_{\geq a}=\{a\}}\\
 &\deleted{\partial\mathcal L=\emptyset}
\end{aligned}
\]

% faint vertical paper watermark
\rotatebox{90}{\textsc{\unsure{ng Extra Stro}}}
